\documentclass[letterpaper, 12pt, titlepage]{article} % Tipo de documento, configurado con: tamaño carta, tamaño de fuente 12pt, inclusión de página de título

% ================================================================
%* PAQUETES BÁSICOS Y CONFIGURACIÓN ESENCIAL
% ================================================================
\usepackage[utf8]{inputenc}         % Codificación de caracteres UTF-8
\usepackage[spanish]{babel}         % Localización en español (guiones, títulos)
\usepackage[margin=1in]{geometry}   % Configuración de márgenes de página

% ================================================================
%* TIPOGRAFÍA Y FORMATO DE TEXTO
% ================================================================
%\usepackage{fontspec}              % Manejo de fuentes con XeLaTeX
%\setmainfont{Myriad Pro} % Fuente principal (Myriad Pro)
\usepackage{helvet}                         % Fuente Sans-Serif (Helvetica)
\renewcommand{\familydefault}{\sfdefault}   % Establece Helvet como fuente principal
\usepackage{setspace}                       % Control de interlineado (1.5, doble espacio)
\usepackage{parskip}                        % Espaciado entre párrafos (en lugar de sangría)
\usepackage{ragged2e}                       % Mejor manejo de justificación de texto

\renewcommand{\thesection}{\Roman{section}} % Numeración romana para secciones
\renewcommand{\thesubsection}{\arabic{section}.\arabic{subsection}} % Subsecciones: 1.1, 1.2...
\renewcommand{\thesubsubsection}{\arabic{section}.\arabic{subsection}.\arabic{subsubsection}} % Subsubsecciones: 1.1.1, 1.1.2...

% ================================================================
%* ELEMENTOS GRÁFICOS Y DISEÑO
% ================================================================
\usepackage{graphicx}               % Manejo de imágenes (inserción y escalado)
\usepackage{xcolor}                 % Sistema de colores personalizados
\usepackage{pdflscape}              % Páginas en orientación horizontal

% ================================================================
%* TABLAS PROFESIONALES
% ================================================================
\usepackage{booktabs}              % Reglas tipográficas para tablas (formato APA)
\usepackage{tabularx}              % Tablas con ancho ajustable
\usepackage{xltabular}             % Tablas largas que ocupan todo el ancho del texto
\usepackage{makecell}              % Celdas personalizadas en tablas
\usepackage{multirow}              % Celdas que abarcan varias filas
% ================================================================
% LISTAS Y ESTRUCTURAS
% ================================================================
\usepackage{enumitem}                   % Personalización avanzada de listas
\setlist{nosep, leftmargin=*, wide=0pt} % Listas ultra compactas

% ================================================================
%* HIPERVÍNCULOS Y METADATOS PDF
% ================================================================
\usepackage[hyphens]{url}
\usepackage{hyperref}             % Hipervínculos y referencias internas
\hypersetup{
    breaklinks=true,                % Configuración de hipervínculos
    colorlinks=true,                % Habilita los enlaces coloreados
    linkcolor=blue,                 % Color de los enlaces internos
    urlcolor=blue,                  % Color de los enlaces externos
    citecolor=blue,                 % Color de las citas
    pdfborder={0 0 0}               % Sin bordes en los enlaces
}

% ================================================================
%* TAREAS PENDIENTES Y ANOTACIONES
% ================================================================
\setlength{\marginparwidth}{2cm} % Soluciona el problema de \todonotes
\usepackage{todonotes} % Paquete principal
\setuptodonotes{
    size=\tiny, % Tamaño de las notas
}
% Comandos personalizados (opcional)
\newcommand{\pendiente}[1]{\todo[color=red!30, inline]{#1}}
\newcommand{\revisar}[1]{\todo[color=yellow!30, inline]{#1}}
\newcommand{\actualizar}[1]{\todo[color=green!30, inline]{#1}}

% ================================================================
%! CONFIGURACIONES FINALES
% ================================================================
\newcommand{\universidad}{Universidad de las Regiones Autónomas de la Costa Caribe Nicaragüense} % Nombre de la universidad
\newcommand{\sede}{Sede de la Universidad} % Nombre de la sede
\newcommand{\carrera}{Nombre de la Carrera} % Nombre de la carrera
\newcommand{\asignatura}{Nombre de la Asignatura} % Nombre de la asignatura
\newcommand{\requisito}{Nombre del Requisito} % Nombre del requisito (asignatura previa)
\newcommand{\codigo}{XXX-XXX} % Código de la asignatura
\newcommand{\docente}{Nombre del Docente, Ing.} % Nombre del docente

% ================================================================
% PAQUETE LOREM IPSUM PARA TEXTO DE EJEMPLO, BORRAR EN PRODUCCIÓN
% ================================================================
\usepackage{lipsum} % Generador de texto de ejemplo

% ================================================================
% INICIO DEL DOCUMENTO (contenido principal)
% ================================================================
\begin{document}

% Se incluyen las secciones del syllabus que se encuentran en el directorio sections
% Portada del documento
% titlepage incluye un \pagebreak (salto de página) automático al final

\begin{titlepage}
    \centering % centrar el contenido de la portada
    % \minipage permite dividir la portada en dos columnas
    \begin{minipage}{0.25\textwidth} % logo en la columna izquierda
        \centering % centrar el logo dentro de la \minipage
        \includegraphics[width=0.6\textwidth]{images/uraccan_logo.png} % inclusión del logo de la universidad
    \end{minipage}%
    \begin{minipage}{0.75\textwidth} % texto en la columna derecha
        \centering % centrar el texto dentro de la \minipage
        \Large{\textbf{\universidad}}\\ % Nombre de la universidad, configurado en PLANTILLA-SILABO-URACCAN.tex
        \Large{\textbf{URACCAN}}\\ % Siglas de la universidad
    \end{minipage}

    \vfill % espacio flexible para centrar el contenido verticalmente
    \large{\textbf{Syllabus}}\\
    \Huge{\textbf{\asignatura}}\\ % Nombre de la asignatura, configurado en PLANTILLA-SILABO-URACCAN.tex

    \vfill % espacio flexible para centrar el contenido verticalmente
    \small{\textbf{\today{}}}\\ % Fecha actual, configurada con \today, que calcula la fecha automáticamente
\end{titlepage}
\section{Datos Generales}
% Para poder usar \footnotes en el título de la sección:
% \SECTION[Título interno de la sección]{Título visible de la sección}
\section[Descriptor de la asignatura]{Descriptor de la asignatura\footnote{Texto íntegro del descriptor, tomado del programa de la asignatura.}}

% Aquí debe poner el texto del descriptor de la asignatura, se encuentra en el programa de la asignatura.

% Texto Lorem Ipsum para el descriptor de la asignatura. Eliminar en producción.
\lipsum[1-4] % Genera cuatro párrafos de texto de ejemplo
\pagebreak
% Para poder usar \footnotes en el título de la sección:
% \SECTION[Título interno de la sección]{Título visible de la sección}
\section[Competencias a Desarrollar u Objetivos Generales]{Competencias a Desarrollar u Objetivos Generales\footnote{Texto íntegro de los objetivos o competencias a desarrollar, tomado del programa de la asignatura.}}

%! Eliminar \PENDIENTE en producción.
\pendiente{Según el enfoque de la asignatura, se pueden definir objetivos generales o compentencias a desarrollar. Pero no los dos a la vez. Revise el programa de la asignatura.}

%! Eliminar \PENDIENTE en producción.
\pendiente{La siguiente enumeración es en caso de que sea una asignatura por competencias. Si es por objetivos, elimine esta tabla y use la enumeración.}

\begin{itemize}[leftmargin=*, itemsep=0.6em]
    \item \lipsum[1][1-3] %! Texto de ejemplo, eliminar en producción
    \item \lipsum[2][1-3] %! Texto de ejemplo, eliminar en producción
    \item \lipsum[3][1-3] %! Texto de ejemplo, eliminar en producción
    \item \lipsum[4][1-3] %! Texto de ejemplo, eliminar en producción
    \item \lipsum[5][1-3] %! Texto de ejemplo, eliminar en producción
    \item \lipsum[6][1-3] %! Texto de ejemplo, eliminar en producción
    \item \lipsum[7][1-3] %! Texto de ejemplo, eliminar en producción
    \item \lipsum[8][1-3] %! Texto de ejemplo, eliminar en producción
    \item \lipsum[9][1-3] %! Texto de ejemplo, eliminar en producción
    \item \lipsum[10][1-3] %! Texto de ejemplo, eliminar en producción
\end{itemize}

%! Eliminar \PENDIENTE en producción.
\pendiente{La siguiente tabla es en caso de que sea una asignatura por objetivos. Si es por competencias, elimine esta tabla y use la primera.}

\small
\begin{xltabular}{\linewidth}{@{}X X X@{}}
    \toprule
    \thead{\small{\textbf{Conceptuales}}} & \thead{\small{\textbf{Procedimentales}}} & \thead{\small{\textbf{Actitudinales}}} \\
    \midrule
    \begin{itemize}[labelsep=2pt, itemsep=0.6em]
        \item \lipsum[1][1-3] %! Texto de ejemplo, eliminar en producción
        \item \lipsum[2][1-3] %! Texto de ejemplo, eliminar en producción
        \item \lipsum[3][1-3] %! Texto de ejemplo, eliminar en producción
        \item \lipsum[4][1-3] %! Texto de ejemplo, eliminar en producción
    \end{itemize} &
    \begin{itemize}[labelsep=2pt, itemsep=0.6em]
        \item \lipsum[1][1-3] %! Texto de ejemplo, eliminar en producción
        \item \lipsum[2][1-3] %! Texto de ejemplo, eliminar en producción
        \item \lipsum[3][1-3] %! Texto de ejemplo, eliminar en producción
        \item \lipsum[4][1-3] %! Texto de ejemplo, eliminar en producción
    \end{itemize} &
    \begin{itemize}[labelsep=2pt, itemsep=0.6em]
        \item \lipsum[1][1-3] %! Texto de ejemplo, eliminar en producción
        \item \lipsum[2][1-3] %! Texto de ejemplo, eliminar en producción
        \item \lipsum[3][1-3] %! Texto de ejemplo, eliminar en producción
        \item \lipsum[4][1-3] %! Texto de ejemplo, eliminar en producción
    \end{itemize} \\
    \bottomrule
\end{xltabular}
\pagebreak
\begin{landscape} % Página en horizontal para la tabla de contenidos
    \section{Plan Temático}

    \doublespacing
    \begin{xltabular}{\linewidth}{@{}c X c c c c c@{}}
        \toprule
        \multirow{2}{*}{\thead{\normalsize{\textbf{No.}}}} & \multirow{2}{*}{\thead{\normalsize{\textbf{Unidad}}}} & \multicolumn{4}{c}{\textbf{Horas}} & \multirow{2}{*}{\thead{\normalsize{\textbf{Créditos}}}}\\ % Celda combinada para las horas teóricas, prácticas, de trabajo independiente y total &
        \cmidrule(lr){3-6}
        && \thead{\normalsize{\textbf{Teóricas}}} & \thead{\normalsize{\textbf{Prácticas}}} & \thead{\normalsize{\textbf{Trabajo Independiente}}} & \thead{\normalsize{\textbf{Total}}} \\
        \midrule
        %! Texto Lorem Ipsum para mostrar el formato de la tabla. Eliminar en producción y escriba cada unidad en cada item.
        %! Las horas están definidas como ejemplo, ajuste según el contenido real de la asignatura.
        I & \lipsum[1][1] & 8 & 8 & 16 & 32 & \multirow{8}{*}{4} \\
        II & \lipsum[2][1] & 8 & 8 & 16 & 32 \\
        III & \lipsum[3][1] & 8 & 8 & 16 & 32 \\
        IV & \lipsum[4][1] & 8 & 8 & 16 & 32 \\
        V & \lipsum[5][1] & 8 & 8 & 16 & 32 \\
        VI & \lipsum[6][1] & 8 & 8 & 16 & 32 \\
        \midrule
        \multicolumn{2}{c}{\textbf{Evaluación}} & 0 & 4 & 8 & 12 \\
        \midrule
        \multicolumn{2}{c}{\textbf{TOTAL}} & 30 & 34 & 128 & 192 & 4\\
        \bottomrule
    \end{xltabular}
\end{landscape}
\section{Mediación Pedagógica}

\subsection*{Unidad I: \lipsum[1][1]} %! Eliminar \lipsum en producción y escribir el título de la unidad

Durante el desarrollo de la asignatura, se aplicará un diseño coherente que articule:
\begin{enumerate}[leftmargin=*, itemsep=0.6em]
    \item \textbf{Metodologías didácticas} (enfoque global).
    \item \textbf{Estrategias didácticas} (plan para unidades o bloques).
    \item \textbf{Técnicas y herramientas} (procedimientos y recursos puntuales).
\end{enumerate}

%* Esta sección describe las metodologías, estrategias y técnicas que se implementarán para facilitar el aprendizaje significativo y activo de los estudiantes, según cada una de las unidades temáticas del curso.

%! Eliminar \pendiente en producción
\pendiente{Convertir de General a Unitario}

\subsection*{Metodologías didácticas}
Son los marcos teórico-prácticos que guían todo el proceso formativo. Para esta asignatura se adoptarán:
\begin{itemize}[leftmargin=*, itemsep=0.6em]
    \item \textbf{Aprendizaje Activo}: paradigma constructivista centrado en la participación y reflexión del estudiante.
    \item \textbf{Aprendizaje Basado en Proyectos (ABP)}: los estudiantes desarrollan proyectos reales o simulados para aplicar conceptos en contextos auténticos.
    \item \textbf{Aprendizaje Basado en Problemas (ABPr)}: se plantean problemas complejos y abiertos que los alumnos investigan y resuelven en grupos.
    \item \textbf{Gamificación}: incorporación de dinámicas de juego (puntos, niveles, recompensas) para motivar y comprometer a los estudiantes.
\end{itemize}

\subsection*{Estrategias didácticas}
Planificaciones adaptadas a cada unidad temática, donde se combinarán técnicas concretas para alcanzar objetivos específicos:
\begin{itemize}[leftmargin=*, itemsep=0.6em]
    \item \textbf{Aula Invertida}: contenidos teóricos en casa (videos, lecturas) y actividades prácticas en clase para resolver dudas y profundizar.
    \item \textbf{Estudio de Casos}: análisis de escenarios reales o hipotéticos, discusión en grupo y toma de decisiones fundamentadas.
    \item \textbf{Simulaciones y Juegos de Rol}: recreación de situaciones profesionales para practicar negociación, resolución de conflictos y toma de decisiones en un entorno controlado.
\end{itemize}

\subsection*{Técnicas y herramientas}
Son procedimientos y recursos específicos que implementan las estrategias:
\begin{itemize}[leftmargin=*, itemsep=0.6em]
    \item \textbf{Uso de Tecnologías Educativas}: plataformas virtuales, aplicaciones y recursos digitales que facilitan la comunicación, la interacción y el aprendizaje autónomo.
    \item \textbf{Retroalimentación Continua e Inmediata}: mecanismos para que los estudiantes reciban comentarios sobre su desempeño en tiempo real, favoreciendo la autoevaluación y el aprendizaje reflexivo.
\end{itemize}
\pagebreak
% Para poder usar \footnotes en el título de la sección:
% \SECTION[Título interno de la sección]{Título visible de la sección}
\section[Evaluación para los Aprendizajes]{Evaluación para los Aprendizajes\footnote{Está basada en el régimen académico de la URACCAN y en el programa de la asignatura.}}

%! Eliminar \pendiente en producción
\pendiente{Debe incluir una descripción de la metodología de evaluación, incluyendo los criterios y procedimientos que se utilizarán para evaluar el aprendizaje de los estudiantes.}
\pendiente{Debe incluir una descripción de los instrumentos de evaluación que se utilizarán, como exámenes, trabajos, proyectos, etc.}
\pendiente{Debe incluir el procedimiento de cálculo de la nota final, incluyendo la nota mínima aprobatoria.}

%! Eliminar \lipsum en producción
\lipsum[1-5]
\pagebreak
\section{Bibliografía}
% Aquí debe poner la bibliografía de la asignatura
%* Use el formato APA para las referencias

%! Eliminar este \pendiente en producción
\pendiente{Agregar bibliografía de la asignatura.}

\subsection*{Referencias Básicas}

%* Ejemplo de referencia a un libro impreso
\hangindent=0.5in
Apellido, N. (Año). \textit{Título del libro}. Editorial.

%* Ejemplo de referencia a un libro en línea
\hangindent=0.5in
Apellido, N. (Año). \textit{Título del libro}. Editorial. \href{https://www.ejemplo.com}{https://www.ejemplo.com}

%* Ejemplo de referencia a un capítulo de un libro
\hangindent=0.5in
Apellido, A. y Apellido, B. (Año). Título del capítulo. En N. Apellido (Ed.), \textit{Título del libro} (pp. xx-xx). Editorial.

%* Ejemplo de página web con contenido estático
\hangindent=0.5in
Apellido, A., Apellido, B., y Apellido, C. (20 de mayo de 2020). \textit{Título del artículo de la página web}. Nombre del sitio web. \href{https://www.ejemplo.com}{https://www.ejemplo.com}

%* Ejemplo de página web con actualizaciones frecuentes
\hangindent=0.5in
Apellido, A., Apellido, B., y Apellido, C. (20 de mayo de 2020). \textit{Título del artículo de la página web}. Nombre del sitio web. Recuperado el día mes año de \href{https://www.ejemplo.com}{https://www.ejemplo.com}

%* Ejemplo de archivo con formato especial adentro de una página web
\hangindent=0.5in
Apellido, A. (03 de agosto de 2020). \textit{Título del archivo} [Archivo Excel]. Nombre del sitio web. \href{https://www.ejemplo.com}{https://www.ejemplo.com}

%* Ejemplo de vídeo de YouTube
\hangindent=0.5in
Nombre del autor. [Nombre de usuario en YouTube] (fecha). \textit{Título del video} [Video]. YouTube. \href{http://youtube.com/url-del-video}{http://youtube.com/url-del-video}

\subsection*{Referencias Complementarias}

%* Ejemplo de referencia a un libro impreso
\hangindent=0.5in
Apellido, N. (Año). \textit{Título del libro}. Editorial.

%* Ejemplo de referencia a un libro en línea
\hangindent=0.5in
Apellido, N. (Año). \textit{Título del libro}. Editorial. \href{https://www.ejemplo.com}{https://www.ejemplo.com}

%* Ejemplo de referencia a un capítulo de un libro
\hangindent=0.5in
Apellido, A. y Apellido, B. (Año). Título del capítulo. En N. Apellido (Ed.), \textit{Título del libro} (pp. xx-xx). Editorial.

%* Ejemplo de página web con contenido estático
\hangindent=0.5in
Apellido, A., Apellido, B., y Apellido, C. (20 de mayo de 2020). \textit{Título del artículo de la página web}. Nombre del sitio web. \href{https://www.ejemplo.com}{https://www.ejemplo.com}

%* Ejemplo de página web con actualizaciones frecuentes
\hangindent=0.5in
Apellido, A., Apellido, B., y Apellido, C. (20 de mayo de 2020). \textit{Título del artículo de la página web}. Nombre del sitio web. Recuperado el día mes año de \href{https://www.ejemplo.com}{https://www.ejemplo.com}

%* Ejemplo de archivo con formato especial adentro de una página web
\hangindent=0.5in
Apellido, A. (03 de agosto de 2020). \textit{Título del archivo} [Archivo Excel]. Nombre del sitio web. \href{https://www.ejemplo.com}{https://www.ejemplo.com}

%* Ejemplo de vídeo de YouTube
\hangindent=0.5in
Nombre del autor. [Nombre de usuario en YouTube] (fecha). \textit{Título del video} [Video]. YouTube. \href{https://www.youtube.com/url-del-video}{https://www.youtube.com/url-del-video}
\pagebreak
\newgeometry{margin=1.5cm} % Ajusta los márgenes para la tabla
\begin{landscape}

    \section{Plan Calendario}

    \small
    \begin{xltabular}{\linewidth}{@{}XXXXXXXX@{}}
        \toprule
        \begin{center}\textbf{Fecha}\end{center} & \begin{center}\textbf{Unidad y Contenido}\end{center} & \begin{center}\textbf{Objetivos}\end{center} & \begin{center}\textbf{Estrategias Metodológicas}\end{center} & \begin{center}\textbf{Trabajo Independiente}\end{center} & \begin{center}\textbf{Recursos y Materiales Didácticos}\end{center} & \begin{center}\textbf{Mecanismos de Evaluación}\end{center} & \begin{center}\textbf{Bibliografía}\end{center} \\
        \midrule
        %\begingroup Sesión 1
        \begin{itemize}[label={}, labelsep=0pt, left=0pt .. 0pt, itemsep=1em] % Fecha de la sesión
            %! Reemplace con la fecha real y detalles de la sesión
            \item díadelasemana, XX de mes de 202X
            \item Sesión \#XX de XX
            \item Presencial o Virtual
        \end{itemize} &
        \vspace{.01em}
        %! Reemplace el número de unidad y título
        \textbf{Unidad X: Título de la Unidad} % Unidad a cubrir
        \vspace{.6em}
        \begin{itemize}[labelsep=2pt, itemsep=0.6em] % Contenidos
            %! Reemplace con los contenidos de la sesión, borrando los comandos \lipsum
            \item \lipsum[1][1]
            \item \lipsum[1][2]
            \item \lipsum[1][3]
            \item \lipsum[1][4]
        \end{itemize} &
                \begin{itemize}[label={}, labelsep=0pt, left=0pt .. 0pt, itemsep=1em] % Objetivos o Resultados de Aprendizaje
            %! Reemplace con los objetivos específicos de la sesión, borrando los comandos \lipsum
            \item \lipsum[2][1-2]
            \item \lipsum[2][1-2]
            \item \lipsum[2][1-2]
        \end{itemize} &
        \begin{itemize}[label={}, labelsep=0pt, left=0pt .. 0pt, itemsep=1em] % Estrategias Metodológicas
            %! Reemplace con las estrategias metodológicas de la sesión
            \item Forma de Organización de la Enseñanza 1
            \item Forma de Organización de la Enseñanza 2
            \item Forma de Organización de la Enseñanza 3
            \item Forma de Organización de la Enseñanza 4
        \end{itemize}                             &
        \begin{itemize}[label={}, labelsep=0pt, left=0pt .. 0pt, itemsep=1em] % Trabajo Independiente
            %! Reemplace con las actividades de trabajo independiente
            \item Trabajo independiente 1
            \item Trabajo independiente 2
            \item Trabajo independiente 3
            \item Trabajo independiente 4
        \end{itemize}                             &
        \begin{itemize}[label={}, labelsep=0pt, left=0pt .. 0pt, itemsep=1em] % Recursos y Materiales Didácticos
            %! Reemplace con los recursos y materiales didácticos de la sesión
            \item Recurso o material 1
            \item Recurso o material 2
            \item Recurso o material 3
            \item Recurso o material 4
        \end{itemize} &
        \begin{itemize}[label={}, labelsep=0pt, left=0pt .. 0pt, itemsep=1em] % Mecanismos de Evaluación
            %! Reemplace con los mecanismos de evaluación de la sesión
            \item Mecanismo de evaluación 1
            \item Mecanismo de evaluación 2
            \item Mecanismo de evaluación 3
            \item Mecanismo de evaluación 4
        \end{itemize} &
        \begin{itemize}[label={}, labelsep=0pt, left=0pt .. 0pt, itemsep=1em] % Bibliografía
            %! Reemplace con la bibliografía de la sesión, borrando los comandos \lipsum
            %! Use el formato APA para las referencias
            \item Bibliografía obligatoria 1
            \item Bibliografía obligatoria 2
            \item Bibliografía obligatoria 3
            \item Bibliografía opcional 1
            \item Bibliografía opcional 2
            \item Bibliografía opcional 3
        \end{itemize}\\
        %\endgroup Sesión 1
        \bottomrule
    \end{xltabular}
\end{landscape}

\end{document}
% ================================================================
% FIN DEL DOCUMENTO
% ================================================================